\documentclass[preview]{standalone}

\usepackage{amsmath}
\usepackage{amssymb}
\usepackage{stellar}
\usepackage{definitions}

\begin{document}

\id{probability-exercises-batch3}
\genpage

\section{Exercises}

\begin{snippetexercise}{probability-exercises-3-ex-22}{}
    Let
    \[
        \varphi(x)=\frac{1}{\sqrt{2\pi}}e^{-\frac{x^2}{2}},
        \qquad x\in\realnumbers.
    \]
    \begin{enumerate}
        \item Show that \(\varphi\) is a \probabilitydensity.
        \item If \(Z\) is a \randomvariable with density \(\varphi\), compute
        \(\expectation[Z]\) and \(\variance(Z)\).
        \item Let \(\mu\in\realnumbers\), let \(\sigma>0\), and define
        \(X=\sigma Z+\mu\). Determine the density, \expectedvalue and
        \variance of \(X\).
    \end{enumerate}
\end{snippetexercise}

\begin{snippetsolution}{probability-exercises-3-ex-22-sol}{}
    \begin{enumerate}
        \item Since \(\varphi(x)\geq0\) for every \(x\in\realnumbers\), it remains
        to show that its integral is \(1\). Let
        \[
            I=\int_{\realnumbers}e^{-\frac{x^2}{2}}\,dx.
        \]
        Then, using polar coordinates in \(\realnumbers^2\),
        \[
            I^2
            =
            \int_{\realnumbers^2}e^{-\frac{x^2+y^2}{2}}\,dx\,dy
            =
            \int_0^{2\pi}\int_0^{+\infty}e^{-\frac{r^2}{2}}r\,dr\,d\theta
            =
            2\pi.
        \]
        Hence \(I=\sqrt{2\pi}\), and
        \[
            \int_{\realnumbers}\varphi(x)\,dx=1.
        \]

        \item By symmetry,
        \[
            \expectation[Z]
            =
            \int_{\realnumbers}z\varphi(z)\,dz
            =
            0.
        \]
        Moreover,
        \[
            \variance(Z)=\expectation[Z^2].
        \]
        Integrating by parts,
        \[
            \expectation[Z^2]
            =
            \frac{2}{\sqrt{2\pi}}
            \int_0^{+\infty}z^2e^{-\frac{z^2}{2}}\,dz
            =
            \frac{2}{\sqrt{2\pi}}
            \int_0^{+\infty}e^{-\frac{z^2}{2}}\,dz
            =
            1.
        \]
        Thus \(\expectation[Z]=0\) and \(\variance(Z)=1\).

        \item Since \(X=\sigma Z+\mu\) and \(\sigma>0\), the change of variables
        \(z=(x-\mu)/\sigma\) gives
        \[
            f_X(x)
            =
            \frac{1}{\sigma}
            \varphi\left(\frac{x-\mu}{\sigma}\right)
            =
            \frac{1}{\sigma\sqrt{2\pi}}
            e^{-\frac{(x-\mu)^2}{2\sigma^2}}.
        \]
        Hence \(X\sim\mathcal N(\mu,\sigma^2)\). Also,
        \[
            \expectation[X]=\sigma\expectation[Z]+\mu=\mu,
            \qquad
            \variance(X)=\sigma^2\variance(Z)=\sigma^2.
        \]
    \end{enumerate}
\end{snippetsolution}

\begin{snippetexercise}{probability-exercises-3-ex-23}{}
    Let
    \[
        \begin{pmatrix}
            X\\
            Y
        \end{pmatrix}
        \sim
        \mathcal N
        \left(
            \begin{pmatrix}
                0\\
                0
            \end{pmatrix},
            Q
        \right),
        \qquad
        Q=
        \begin{pmatrix}
            a & b\\
            b & c
        \end{pmatrix},
    \]
    where \(Q\) is a \positivedefinitematrix. Define
    \[
        V=X,
        \qquad
        Z=X+Y.
    \]
    Determine the law of \((V,Z)\) and the law of \(X+Y\). What happens if
    \(X\) and \(Y\) are independent?
\end{snippetexercise}

\begin{snippetsolution}{probability-exercises-3-ex-23-sol}{}
    We have
    \[
        \begin{pmatrix}
            V\\
            Z
        \end{pmatrix}
        =
        A
        \begin{pmatrix}
            X\\
            Y
        \end{pmatrix},
        \qquad
        A=
        \begin{pmatrix}
            1 & 0\\
            1 & 1
        \end{pmatrix}.
    \]
    By stability of \gaussianvector[Gaussian vectors] under invertible affine
    transformations,
    \[
        \begin{pmatrix}
            V\\
            Z
        \end{pmatrix}
        \sim
        \mathcal N(0,\Sigma),
        \qquad
        \Sigma=AQA^T.
    \]
    Therefore
    \[
        \Sigma
        =
        \begin{pmatrix}
            1 & 0\\
            1 & 1
        \end{pmatrix}
        \begin{pmatrix}
            a & b\\
            b & c
        \end{pmatrix}
        \begin{pmatrix}
            1 & 1\\
            0 & 1
        \end{pmatrix}
        =
        \begin{pmatrix}
            a & a+b\\
            a+b & a+c+2b
        \end{pmatrix}.
    \]
    The second marginal is therefore
    \[
        X+Y=Z\sim\mathcal N(0,a+c+2b).
    \]
    Since \(a=\variance(X)\), \(c=\variance(Y)\), and \(b=\covariance(X,Y)\),
    this can be written as
    \[
        X+Y
        \sim
        \mathcal N\left(
            0,
            \variance(X)+\variance(Y)+2\covariance(X,Y)
        \right).
    \]
    If \(X\) and \(Y\) are independent, then \(b=\covariance(X,Y)=0\), so
    \[
        X+Y\sim\mathcal N(0,\variance(X)+\variance(Y)).
    \]
\end{snippetsolution}

\begin{snippetexercise}{probability-exercises-3-ex-24}{}
    Let \((X,Y)\) be an \absolutelycontinuousrandomvector in
    \(\realnumbers^2\) with joint density
    \[
        f_{X,Y}(x,y)
        =
        \frac{x}{4}e^{-\frac{x^2}{4}}y\,\indicator_D(x,y),
        \qquad
        D=\{(x,y)\in\realnumbers^2\suchthat 0<y<x\}.
    \]
    \begin{enumerate}
        \item Compute the marginal densities of \(X\) and \(Y\).
        \item Determine whether \(X\) and \(Y\) are independent.
        \item Let
        \[
            U=X,
            \qquad
            V=\frac{X}{Y}.
        \]
        Determine whether \(U\) and \(V\) are independent.
    \end{enumerate}
\end{snippetexercise}

\begin{snippetsolution}{probability-exercises-3-ex-24-sol}{}
    \begin{enumerate}
        \item For \(x>0\),
        \[
            f_X(x)
            =
            \int_0^x\frac{x}{4}e^{-\frac{x^2}{4}}y\,dy
            =
            \frac{x^3}{8}e^{-\frac{x^2}{4}}.
        \]
        Hence
        \[
            f_X(x)=\frac{x^3}{8}e^{-\frac{x^2}{4}}\indicator_{(0,+\infty)}(x).
        \]
        For \(y>0\),
        \[
            f_Y(y)
            =
            \int_y^{+\infty}\frac{x}{4}e^{-\frac{x^2}{4}}y\,dx
            =
            \frac{y}{2}e^{-\frac{y^2}{4}}.
        \]
        Hence
        \[
            f_Y(y)=\frac{y}{2}e^{-\frac{y^2}{4}}\indicator_{(0,+\infty)}(y).
        \]

        \item The variables \(X\) and \(Y\) are not independent. Indeed, if
        \(0<x<y\), then \(f_{X,Y}(x,y)=0\), while \(f_X(x)>0\) and
        \(f_Y(y)>0\). Thus the joint density cannot factor as the product of
        the two marginal densities.

        \item Consider the map
        \[
            \varphi(x,y)=\left(x,\frac{x}{y}\right).
        \]
        On \(D\), its image is
        \[
            \varphi(D)=(0,+\infty)\cartesianprod(1,+\infty),
        \]
        and its inverse is
        \[
            \varphi^{-1}(u,v)=\left(u,\frac{u}{v}\right).
        \]
        The absolute value of the Jacobian determinant of \(\varphi^{-1}\) is
        \[
            \left|\det J\varphi^{-1}(u,v)\right|=\frac{u}{v^2}.
        \]
        Therefore
        \[
            f_{U,V}(u,v)
            =
            \frac{u^3}{4v^3}e^{-\frac{u^2}{4}}
            \indicator_{(0,+\infty)}(u)\indicator_{(1,+\infty)}(v).
        \]
        This factorizes as
        \[
            f_{U,V}(u,v)
            =
            \left(
                \frac{u^3}{8}e^{-\frac{u^2}{4}}
                \indicator_{(0,+\infty)}(u)
            \right)
            \left(
                \frac{2}{v^3}\indicator_{(1,+\infty)}(v)
            \right).
        \]
        Both factors are densities, so \(U\) and \(V\) are independent.
    \end{enumerate}
\end{snippetsolution}

\begin{snippetexercise}{probability-exercises-3-ex-25}{}
    Let \((X,Y)\) be uniformly distributed on the square \([-1,1]^2\), so that
    \[
        f_{X,Y}(x,y)=\frac{1}{4}\indicator_{[-1,1]^2}(x,y).
    \]
    Let
    \[
        G=
        \frac{1}{\sqrt2}
        \begin{pmatrix}
            1 & -1\\
            1 & 1
        \end{pmatrix}
    \]
    and define
    \[
        \begin{pmatrix}
            U\\
            V
        \end{pmatrix}
        =
        G
        \begin{pmatrix}
            X\\
            Y
        \end{pmatrix}.
    \]
    Determine the joint density of \((U,V)\), and show that \(U\) and \(V\)
    are not independent.
\end{snippetexercise}

\begin{snippetsolution}{probability-exercises-3-ex-25-sol}{}
    The matrix \(G\) is orthogonal, so \(G^{-1}=G^T\) and
    \(|\det G|=1\). Hence
    \[
        G^{-1}
        \begin{pmatrix}
            u\\
            v
        \end{pmatrix}
        =
        \frac{1}{\sqrt2}
        \begin{pmatrix}
            u+v\\
            v-u
        \end{pmatrix}.
    \]
    By the affine density transformation,
    \[
        f_{U,V}(u,v)
        =
        \frac{1}{4}
        \indicator_D(u,v),
    \]
    where
    \[
        D=
        \left\{
            (u,v)\in\realnumbers^2
            \suchthat
            |u+v|\leq\sqrt2,\ |v-u|\leq\sqrt2
        \right\}.
    \]
    This is the square \([-1,1]^2\) rotated by \(45^\circ\).

    At the point \((0,0)\), the joint density is
    \[
        f_{U,V}(0,0)=\frac{1}{4}.
    \]
    The marginal densities at \(0\) are
    \[
        f_U(0)
        =
        \int_{-\sqrt2}^{\sqrt2}\frac{1}{4}\,dv
        =
        \frac{\sqrt2}{2},
        \qquad
        f_V(0)
        =
        \int_{-\sqrt2}^{\sqrt2}\frac{1}{4}\,du
        =
        \frac{\sqrt2}{2}.
    \]
    Therefore
    \[
        f_U(0)f_V(0)=\frac{1}{2}\neq\frac{1}{4}=f_{U,V}(0,0).
    \]
    Thus the joint density does not factor into the product of the marginals,
    and \(U\) and \(V\) are not independent.
\end{snippetsolution}

\end{document}
